\documentclass[10pt,a4paper]{amsart}
\usepackage[margin=19mm]{geometry}
\usepackage{amsmath,amssymb,mathtools}
\usepackage[scr=rsfs,cal=euler]{mathalfa}
\usepackage{xfrac}
\usepackage[unicode,hidelinks,bookmarks=false]{hyperref}
\newcommand{\ie}{i.e.\ }
\newcommand{\cf}{cf.\ }
\newcommand{\resp}{respectively\ }
\newcommand{\Subsection}{\subsection}
\providecommand{\nobreakdash}{\nobreak\hbox{-}}
\setlength{\emergencystretch}{2em}
\multlinegap0pt
\usepackage{amssymb}
\usepackage{bm}
\usepackage{braket}
\usepackage{dsfont}
\usepackage{csquotes}
\newtheorem{definition}{Definition}
\newtheorem{proposition}{Proposition}
\newtheorem{theorem}{Theorem}
\newcommand{\T}{\ms T}
\newcommand{\bP}{{\mathbf P}}
\newcommand{\bX}{\partial_\infty X}
\newcommand{\defeq}{\coloneqq}
\newcommand{\im}{\operatorname{Im}}
\newcommand{\mc}{\mathcal}
\newcommand{\ms}{\mathsf}
\newcommand{\seq}[1]{ \{{#1}_p\}_{p\in\mathbb N}}
\title{The topological Ax--Lindemann theorem\\ for the complex hyperbolic space}
\author{Fran\c{c}ois Labourie, Shahar Mozes}
\date{Source: arXiv:2609.06440v1 (CC BY 4.0)}
\begin{document}
\maketitle

\noindent\emph{Source and licence.} The topological Ax--Lindemann theorem\\ for the complex hyperbolic space. Version arXiv:2609.06440v1 (CC BY 4.0), distributed by arXiv under the Creative Commons Attribution 4.0 International licence, http://creativecommons.org/licenses/by/4.0/, as recorded in the arXiv licence metadata for this exact version and shown on the abstract page https://arxiv.org/abs/2609.06440v1. Full licence notice: \texttt{licenses/arxiv-2609.06440-CC-BY-4.0.txt}.\par\smallskip

\noindent\textit{Editorial scope.} Counted unit: the article's theorem theorem (source0.tex lines 594--600). The article's complete original proof of it is delivered with it (source0.tex lines 602--651, one \texttt{proof} environment of 50 lines). No mathematics is written, repaired or re-derived: the statement and the proof are the article's own lines, in the article's own notation, and no retained line is substituted or re-flowed.  Retained uncounted as the source-local context the proof needs: lines 295--302; lines 435--439; lines 543--548; lines 570--573.  Nothing was omitted from any retained span, so the package carries no omission record.  No retained line contains a citation, so the wrapper carries no bibliography and no bibliography entry.  No retained line contains a figure environment, an image-inclusion command or drawing code, so no image is delivered and the package declares no asset dependency.  Editorial formatting only, in addition: an AMS \texttt{amsart} preamble that restates the article's own declarations and macros used by the retained lines, namely the environments \texttt{definition}, \texttt{proposition}, \texttt{theorem} on separate counters, together with the standard packages those lines need.  The retained environments print in this document as Definition 1, Proposition 1, Theorem 1 in order of appearance, whereas the article prints the same items with the numbering of its own full document.  The complete native source of the article as distributed in the arXiv e-print for this exact version is bundled unchanged as \texttt{sources/209536/source0.tex}.
\begin{definition}[\sc Hypothesis $(*)$]\label{sec:hypo*}
	We say $(Z,\phi)$  {\em satisfies hypothesis $(*)$} if  
\begin{enumerate}
	\item the set of  points $x$ in $Z$  where the tangent map  $\T_x\phi$ is injective is non empty (hence open and dense).	\item $\phi(Z)\cap X$ is non empty,
	\item $\phi(\partial Z)$ is a subset of ${\bf P}(E)\setminus X$.
\end{enumerate} 

\end{definition} 

\begin{proposition}\label{pro:regnotin}
	Assume $\phi(Z)$ is not included in a complex projective space $M$ whose projection on $\Gamma\backslash X$ is closed.
	
	Then there exists a regular element $z$ such that $\phi(z)$ does not belong to any  $N$ in $\mathcal N$.
\end{proposition}

\begin{proposition}
	\label{pro:lem1} Let $(Z,\phi)$ satisfying hypothesis $(*)$ as in section \ref{sec:hypo*}, $y$ a regular point (identified with $\phi(y)$) and $L$ the complex complex affine hypersurface through $y$, such that $\T_yL=\T_y Z$. Then
$$
W_y(L)\subset W_y(Z)\ .	
$$
\end{proposition}

\begin{proposition}\label{pro:wldense}
	Assume that $L$ is a complex line  transverse to $\bX$. Assume that $y$ is a  point in  in $\partial_\infty L$ and $y$  does not belong to  $\partial_\infty N$ for any $N$ in $\mc N$.
	Then $W_y(L)$ is dense.
\end{proposition}

\begin{theorem}
	Let $\Gamma$ a uniform lattice in the group of biholomorphisms of $X$.
	
	Let $Z$ be a connected complex manifold with boundary, $\phi$ a holomorphic map with values in ${\bf P}(E)$ which is an immersion somewhere. Assume that $\phi(Z)$ intersects $X$ and  $\phi(\partial Z)$ lies in ${\bf P}(E)\setminus X$. 
	
Let $\pi$ be the projection from $X$ to $\Gamma\backslash X$, then  $$\pi(\phi(Z)\cap X)\ ,$$ is dense (for the usual topology) in a totally geodesic subspace of  $\Gamma\backslash X$.
\end{theorem}

\begin{proof} 
We can safely assume that $\phi(Z)$ is not included in any $N$ in $\mc N$. Otherwise, we would work on $N_0$ which is the complex hyperbolic subspace of the lowest possible dimension projecting to a closed totaly geodesic subspace in $\Gamma\backslash X$.

By proposition \ref{pro:regnotin}, there exists a regular element $y$ in $Z$, such that $z\defeq\phi(y)$ does not belong to 
 $$
 \bigcup_{N\in\mc N}\partial_\infty N\ .$$  
 \vskip 0,5 truecm
\noindent{\sc First Step: dimensional reduction.}
Our first step is to reduce to the case  when $Z$ has complex dimension 1. 
 Let $p$ be the dimension of $Z$ and assume $p>1$ for the moment. 
 Let  $\mc F$ the set of complex projective spaces of codimension $p-1$ and intersecting $X$. 
For every $F$ in $\mc F$, let $Z_F\defeq \phi^{-1}(F)$.

There is an non empty open set $\mc F_0$ in $\mc F$, such that for any $F$ in $\mc F_0$, 
\begin{enumerate}
	\item $\phi$ is transverse to $F$,
	\item $\phi(Z_F)$ intersects $X$.
\end{enumerate}
and thus $Z_F$ is a curve with $\partial Z_F\subset \partial Z$. In particular $(Z_F,\phi\vert_{Z_F})$ satisfy hypothesis $(*)$ of definition \ref{sec:hypo*} for all $F$ in $\mc F_0$.   

For every $N$ in $\mc N$, let us consider the following subset of $\mc F_0$.
$$
{\mc F}(N)=\{F\in \mc F_0\mid \phi(Z_F)\subset N\}.
$$ 
Observe now for any $N$ in $\mc N_0$, the subset  ${\mc F}(N)$ is closed: indeed let $\seq{F}$ be a sequence in  ${\mc F}(N)$ converging to $F$ in $\mc F_0$, then for any point $z$ in $Z_F$ there is a sequence of point $\seq{z}$ with $z_p$ in $Z_{F_p}$ converging to $z$. Hence $z$ belongs to $N$, and since this is true for all $z$ in $Z_F$, $Z_F$ is a subset of $N$.

Furthermore assume that   
${\mc F}(N)$ contains an open set $O$, then the following subset of $Z$
$$
U\defeq \bigcup_{F\in O}Z_F 
$$
contains an open set. Since $U\subset N$, by analytic continuation this implies that $\phi(Z)\subset N$ hence that $N=\bP(E)$. 

If follows that by Baire Theorem
$$
\bigcup_{N\in\mc N}{\mc F}(N)\ ,
$$
 is a Baire meagre set. In particular, its complementary $O_{\mc N}$ in $\mc F_0$ is nonmeagre and thus non empty. For $F$ in  $O_{\mc N}$ we have by construction that $(Z_F,\phi\vert_{Z_F})$ is a curve satisfying hypothesis $(*)$ and not included in a 
 $$
\bigcup_{N\in\mc N} N\ .
 $$
It is enough to prove that $\pi(\phi(Z_F))$ is dense. In other words we have reduced to the case of curves. 
\vskip 0,5 truecm
\noindent{\sc Final Step:} We now can assume $Z$ is a curve. Let $z$ be a point in $\bX$ and not in 
$$
\bigcup_{N\in\mc N}\partial_\infty N\ .
$$ 
Let $L$ be any complex line through $z$ transverse to $\bX$. By proposition \ref{pro:wldense}, the limiting set $W_z(L)$ is dense. Take now $z=\phi(y)$ where $y$ is regular in $Z$. Let now $L=\im(\T_y\phi)$ which is transverse at $y$ to $\bX$ since $y$ is regular. Then by proposition \ref{pro:lem1},
$W_z(L)\subset W_y(Z)$. Thus $W_y(Z)$ is dense. We conclude  by noticing that, by construction,  $p(W_y(Z))$ is included in the closure of  $\pi(\phi(Z)\cap X)$.
	\end{proof}

\end{document}
